%============================================================
\chapter{Aplicación al Proyecto Fin de Curso}\label{cap:actividades}
%============================================================

Este capítulo traslada lo aprendido con BiblioUTEQ al Proyecto Fin de Curso (PFC) de cada equipo. Las actividades siguen el orden de la Unidad 5: primero se revisa la arquitectura, luego se refactoriza y por último se verifica la calidad.

\section{Señales de que las capas se están erosionando}

Antes de modificar el PFC, revise si presenta alguno de estos problemas. Cada uno es una forma concreta de la erosión de la arquitectura descrita en el capítulo \ref{cap:fundamentos} \cite{li2022understanding}.

\begin{xltabular}{\linewidth}{@{}L{4.2cm}YL{4.3cm}@{}}
	\caption{Señales de erosión de la arquitectura en capas}\label{tab:erosion}\\
	\toprule
	\rowcolor{uteqBlue}\hdr{Señal} & \hdr{Por qué es un problema} & \hdr{Cómo se corrige} \\
	\midrule
	\endfirsthead
	\toprule
	\rowcolor{uteqBlue}\hdr{Señal} & \hdr{Por qué es un problema} & \hdr{Cómo se corrige} \\
	\midrule
	\endhead
	El controlador usa un repositorio directamente & Se salta la capa de negocio; las reglas quedan dispersas & Mover la lógica a un servicio \\
	\addlinespace
	El API devuelve entidades JPA & Expone datos internos y ata el contrato al esquema & DTO de respuesta (capítulo \ref{cap:servicios}) \\
	\addlinespace
	Reglas del negocio en el controlador o en la vista & La misma regla se repite y se contradice entre clientes & Llevarla al dominio o al servicio \\
	\addlinespace
	\texttt{@Transactional} en el controlador & La transacción incluye trabajo de presentación y dura más de lo necesario & Delimitar la transacción en el servicio \\
	\addlinespace
	Servicios que solo reenvían al repositorio & Capas que no aportan nada; las solicitudes atraviesan capas vacías \cite{richards2025fundamentals} & Aceptarlo si es poco frecuente; si es la mayoría, revisar el diseño \\
	\addlinespace
	Llamadas remotas sin tiempo de espera & Una falla en otro nivel bloquea al que llama & Configurar tiempos de espera (capítulo \ref{cap:presentacion}) \\
	\addlinespace
	Excepciones técnicas que llegan al cliente & El cliente recibe trazas internas y no puede reaccionar & Manejador global con Problem Details \\
	\bottomrule
\end{xltabular}

\section{Actividades}

\begin{practicaBox}[title={Actividad 11.1: diagnóstico de capas del PFC}]
	Para el microservicio principal de su PFC, dibuje sus capas y las dependencias reales entre sus paquetes (puede usar el diagrama de dependencias de IntelliJ IDEA). Marque con rojo cada dependencia que viole la regla de la sección \ref{sec:regla} y cada señal de la tabla \ref{tab:erosion} que encuentre. Entregue el diagrama y la lista priorizada de correcciones en el documento del PFC.
\end{practicaBox}

\begin{practicaBox}[title={Actividad 11.2: reglas de capas como pruebas}]
	Agregue la biblioteca ArchUnit a las pruebas del PFC y escriba al menos tres reglas que fallen si alguien viola la arquitectura: (a) las clases del paquete de negocio no dependen del paquete web; (b) los controladores no dependen de los repositorios; (c) las entidades no se usan en las firmas públicas de los controladores. Ejecútelas en la integración continua: a partir de ese momento, la erosión se detecta automáticamente.
\end{practicaBox}

\begin{practicaBox}[title={Actividad 11.3: refactorización guiada}]
	Corrija las tres violaciones de mayor prioridad de la actividad 11.1. Por cada una, haga un \emph{commit} independiente cuyo mensaje explique la regla que se restablece, y verifique que las pruebas existentes siguen pasando antes y después del cambio.
\end{practicaBox}

\begin{practicaBox}[title={Actividad 11.4: matriz de calidad del PFC}]
	Construya para su PFC una tabla como la tabla \ref{tab:iso}: para cada característica de ISO/IEC 25010:2023 que su equipo priorice, indique la decisión de diseño que la favorece y la evidencia verificable en el repositorio (prueba, medición, configuración o captura). No incluya características sin evidencia.
\end{practicaBox}

\section{Lista de verificación de una aplicación distribuida en capas}

\begin{okBox}[title={Antes de entregar, compruebe que su sistema cumple}]
	\begin{itemize}
		\item[$\square$] Cada capa usa solo a la inmediatamente inferior; ninguna clase de negocio importa clases web.
		\item[$\square$] Los controladores no contienen reglas del negocio y no usan repositorios.
		\item[$\square$] El API usa DTO, valida su entrada y responde errores con Problem Details.
		\item[$\square$] Las transacciones se delimitan en la capa de negocio.
		\item[$\square$] Los datos que pueden modificarse a la vez tienen control de concurrencia y una prueba que lo demuestra.
		\item[$\square$] El esquema de la base de datos se crea con migraciones versionadas.
		\item[$\square$] Toda llamada entre niveles tiene tiempo de espera y un manejo controlado de la falla.
		\item[$\square$] La configuración que cambia entre ambientes (URL, usuarios, claves) se lee de variables de entorno.
		\item[$\square$] El sistema completo se levanta con \texttt{docker compose up} desde un clon limpio.
		\item[$\square$] Existen pruebas por capa y de integración, y la integración continua está en verde.
		\item[$\square$] El servicio publica su salud y sus métricas.
	\end{itemize}
\end{okBox}
